\begin{recipe}[
    preparationtime = {1h 30min},
    bakingtime = {50min},
    bakingtemperature = {\protect\bakingtemperature{
        topbottomheat=\unit[200]{\textcelcius}
    }},
    portion = {\portion{8}},
    calory = {780kcal | 14g Protein | 52g KH | 55g Fett},
    price = {2,62€},
    rating = {4},
    source = {\protect\recipesource{https://cooking.nytimes.com/recipes/1020596-vegetarian-mushroom-wellington}{NYTimes Cooking: Vegetarian Mushroom Wellington}}
]{Mushroom Wellington}

\graph{% pictures
        small=pic/hauptgericht/07_veggi-beef-wellington-small,
        big=pic/hauptgericht/07_veggi-beef-wellington-big
    }

\ingredients{
    \textbf{Pilze} & \\
    4 & Portobello Pilze\\
    900g & Gemischte Pilze\\
    200ml & Olivenöl (gesamt)\\
    4 & Schalotten\\
    6 & Knoblauchzehen\\
    2 EL & Rosmarin\\
    1 TL & Thymian\\
    80ml & Portwein\\
    2 EL & Sojasauce\\
    100g & Walnüsse\\
    \\
    \textbf{Zwiebeln} & \\
    3 EL & Butter\\
    2 & Gelbe Zwiebeln\\
    240ml & Apfelsaft\\
    1 EL & Balsamico (optional)\\
    \\
    \textbf{Teig} & \\
    400g & Blätterteig\\
    1 & Ei\\
    Mehl & zum Bestäuben\\
    \multicolumn{2}{l}{\textit{\small Fortsetzung auf nächster Seite}}
}
\ingredientsB{
    \textbf{Sauce} & \\
    360ml & Portwein\\
    360ml & Gemüsebrühe\\
    3 EL & Butter\\
    1 & Schalotte\\
    2 & Knoblauchzehen
}

\preparation{
    \step%
    Portobellos putzen, mit Öl bepinseln und je 5min pro Seite anbraten. Abkühlen lassen.
    
    \step%
    Gemischte Pilze fein hacken.
    
    \step%
    In Olivenöl Pilze, Schalotten, Knoblauch und Kräuter 10min kräftig anbraten.
    
    \step%
    Port und Sojasauce zugeben und einkochen lassen. Walnüsse unterrühren und abkühlen.
    
    \step%
    Zwiebeln mit Butter, Zucker und Gewürzen karamellisieren, Apfelsaft zugeben und einkochen.
    
    \step%
    Ofen auf 200°C vorheizen. Blätterteig ausrollen.
    
    \step%
    Hälfte der Pilzmasse mittig verteilen, Zwiebeln darauf, Portobellos darauflegen.
    
    \step%
    Restliche Pilzmasse darüber geben und zu einer kompakten Form drücken.
    
    \step%
    Blätterteig darüber schlagen, mit Ei bestreichen und verschließen.
    
    \step%
    45--50min backen bis goldbraun.
    
    \step%
    Für die Reduktion Port und Brühe einkochen, Butter einrühren.
}

\suggestion[Festlich]{Perfekt als vegetarisches Weihnachts-Hauptgericht mit Rotkohl und Kartoffelklößen.}

\hint{Die Pilzmasse muss vollständig abgekühlt sein, bevor sie in den Blätterteig kommt, sonst wird der Teig weich.}

\end{recipe}